% with lualatex
\documentclass[12pt,french]{article}
\usepackage[nonamssymb]{ProfCollege}
\usepackage{ProfMaquette}
\input{Christ7}
\usepackage[warnings-off={mathtools-colon,mathtools-overbracket}]{unicode-math}
\usepackage{fourier-otf}
\usepackage[a4paper,margin=1.5cm,nohead,includefoot]{geometry}
\setlength{\parindent}{0pt}
\pagestyle{empty}
\usepackage{mflogo}
\usepackage{babel}

\def\MPclipOne#1#2#3{
  \begin{mplibcode}
    vardef paired (expr d) =
    if pair d : d else : (d,d) fi
    enddef ;
    
    primarydef p xyscaled q =
    begingroup ; save qq ; pair qq ; qq = paired(q) ;
    ( p
    if xpart qq<>0 : xscaled (xpart qq) fi
    if ypart qq<>0 : yscaled (ypart qq) fi )
    endgroup
    enddef ;
    
    path fullsquare;
    fullsquare := unitsquare shifted - center unitsquare ;
    
    color myyellow;
    myyellow = (242/255,226/255,149/255);
    
    path p; p := fullsquare xyscaled (#1,#2);% squeezed #3;

    %fill (point(0) of p)--(point(1) of p)--(point(2) of p)--(point(3) of p)--cycle;
    pair pp;
    pp=0.5[point(0) of p,point(2) of p];
    height=2pt;
    width=2pt;
    p:=(point(0) of p)--(point(1) of p)--(point(2) of p);

    p:=(1/3[point(0) of p,point(1) of p])--(point(1) of p)--(2/3[point(1) of p,point(2) of p]);% withpen pencircle scaled 2bp;
    z1=point(0) of p;
    x2=x12=xpart(point(length p/6) of p);
    x3=x11=xpart(point(2*length p/6) of p);
    x4=xpart(point(3*length p/6) of p)+width;
    x5=x4-(width/3);
    x6=x4-(2*width/3);
    x7=xpart(point(length p) of p);
    x8=x7-(width/3);
    x9=x7-(2*width/3);
    x10=x7-width;
    y2=ypart(point(length p/6) of p)-(height/3);
    y3=ypart(point(2*length p/6) of p)-(2*height/3);
    y4=ypart(point(3*length p/6) of p)-height;
    y12=ypart(point(length p/6) of p)+(height/3);
    y11=ypart(point(2*length p/6) of p)+(2*height/3);
    y10=ypart(point(3*length p/6) of p)+height;
    y5=y9=ypart(point(4*length p/6) of p);
    y6=y8=ypart(point(5*length p/6) of p);
    y7=ypart(point(length p) of p);
    fill z1--z2--cycle--z3--z4--z5--z6--z7--z8--z9--z10--z11--z12--cycle;% withcolor 0.85myyellow;
    fill (z1--z2--cycle--z3--z4--z5--z6--z7--z8--z9--z10--z11--z12--cycle) rotatedabout(pp,180);% withcolor 0.85myyellow;
  \end{mplibcode}
}

\newsavebox{\tacochapterbox}

\newcommand{\Tacoo}[1]{%
  \begin{lrbox}{\tacochapterbox}
    #1%
  \end{lrbox}
  \ooalign{%
    \MPclipOne{\mpdim{\wd\tacochapterbox+1pc}}% width
    {\mpdim{\ht\tacochapterbox+\dp\tacochapterbox+1pc}}% height
    {0pt}% thickness of the curve
    \cr\kern0.5pc\raisebox{\ht\tacochapterbox\relax}{\usebox\tacochapterbox}%
  }
}

\NewDocumentEnvironment{CadreMP}{+b}{%
  \begin{lrbox}{\tacochapterbox}
    \begin{minipage}{0.97\linewidth}
      #1%
    \end{minipage}
  \end{lrbox}
  \par
  \ooalign{%
    \MPclipOne{\mpdim{\wd\tacochapterbox+1pc}}% width
    {\mpdim{\ht\tacochapterbox+\dp\tacochapterbox+1pc}}% height
    {0pt}% thickness of the curve
    \cr\kern0.5pc\raisebox{\ht\tacochapterbox\relax}{\usebox\tacochapterbox}%
  }%
}{}%

\usepackage{pdfpages}

\title{Factoriser la production de documents}
\author{C.Poulain}
\date{\PfMfiledate}

\usepackage[colorlinks=true,linkcolor=purple]{hyperref}

\pagestyle{empty}

\usepackage{listings}
\lstset{
    language     = [LaTeX]TeX,
    basicstyle   = \ttfamily,%
    breaklines = true,
    commentstyle = \footnotesize\slshape\color{gray},
    emphstyle=\color{purple},
    columns=fullflexible,%
    frame=tb,%
    texcsstyle=*\color{black},%pour colorer la contre-oblique
    classoffset=5,
    texcsstyle=*\color{blue},%pour colorer la contre-oblique
    moretexcs={brm,TikzDM,TikzDS,TikzIE,TikzFiche,Competences,TikzParcours,TikzPdT,PfMCompNA,PfMCompECA,PfMCompA,AfficheParcours},
    emph={Maquette,exercice,Solution,Reponse,Indice},%
    classoffset=1,
    keywords={DM,DS,IE,Fiche,CorrigeFin,CorrigeApres,PdT,Parcours,ParcoursPerso},
    keywordstyle=\color{OliveDrab},
    classoffset=2,
    keywords={Theme,Niveau,Classe,Date,Calculatrice,Code,Numero,Sujet,Nom,NumSujet,Fichier,Type},
    keywordstyle=\color{BlueViolet},%DarkMagenta},
    classoffset=3,
    keywords={Pouce,Direct,PasCorrige,BaremeDetaille,BaremeTotal,Source,Titre,Oral,Calculatrice,Competence,Trajet},
    keywordstyle=\color{DarkGreen},
%    delim        = **[s][\color{purple}]{$}{$},
%    moredelim    = **[s][\color{purple}]{\\[}{\\]},
%    moredelim    = **[is][\color{black}]{\\Z}{\\Z},
    literate=*{[}{{\textcolor{orange}{[}}}{1}
    {]}{{\textcolor{orange}{]}}}{1}
    {\{}{\textcolor{amber}{\{}}{1}
    {\}}{\textcolor{amber}{\}}}{1}
    {\&}{\textcolor{red}{\&}}{1}
    {\\[}{{\textcolor{purple}{\textbackslash[}}}{2}
    {\\]}{{\textcolor{purple}{\textbackslash]}}}{2}
    {$}{{\textcolor{purple}{\$}}}{1}%$
    {_}{{\textcolor{purple}{\_}}}{1}%
    {^}{{\textcolor{purple}{\^{}}}}{1}%,
}

% "Rehook" delimiters char table.
\makeatletter
\lst@AddToHook{SelectCharTable}{\lst@DefDelims}
\makeatother

\usepackage{textcomp}%pour les < de listings
\usepackage{enumerate}
\usepackage{hhline}
\usepackage{pifont}

\definecolor{mygray}{RGB}{245,245,245}%pour le background de listing
\definecolor{drab}{rgb}{0.59, 0.44, 0.09}
\definecolor{amber}{rgb}{1.0, 0.75, 0.0}
\definecolor{chocolate}{rgb}{0.82, 0.41, 0.12}
\colorlet{amber}{chocolate}

\newtcblisting{Codes}[3][]{%
  top=0mm,bottom=0mm,left=2mm,right=2mm,middle=0mm,%
  colback=white,%
  colframe=white!75!black,%
  every listing line={#3},%
  listing options={%
    frame=,%
    backgroundcolor=,%
  },%
  righthand width=#2\linewidth,%
  #1%
}%

\usepackage{menukeys}
\newcommand\metamk[1]{\textlangle#1\textrangle}

\newcommand\Cle[1]{%
  {\sffamily\textlangle #1\textrangle}%
}%

\newcommand{\Defaut}[1]{%
\hfill valeur par défaut : {\sffamily #1}\par%
}%

\usepackage{parskip}

\newenvironment{Description}
               {\list{}{\labelwidth0pt \leftmargin0pt \itemindent-\leftmargin
                        \let\makelabel\Descriptionlabel}}
               {\endlist}
\newcommand*\Descriptionlabel[1]{\hspace\labelsep
                                \normalfont\bfseries #1}

\parindent0pt

\newcommand\TIKZ{Ti\textit{k}Z}

\renewcommand\TikzDM{%
  \begin{tcolorbox}[frame hidden,colback=white,enhanced,%
    borderline north={3pt}{0pt}{gray!85},
    borderline north={2pt}{0.5pt}{gray!15},
    borderline south={3pt}{0pt}{gray!85},
    borderline south={2pt}{0.5pt}{gray!15},
    ]%
    \sffamily Devoir en temps libre \no\useKV[DM]{Numero}\hfill\useKV[DM]{Niveau}\ieme{} \useKV[DM]{Classe}%
    \par{\tiny\jobname}\hfill{\scriptsize Pour le \useKV[DM]{Date}}%
  \end{tcolorbox}%
}
\renewcommand\TikzDS{%
  \begin{tcolorbox}[colback=gray!5,%
    enhanced,%
    overlay unbroken and first={%
      \node[yshift=1em] at (frame.south) {\scriptsize\sffamily-- Calculatrice \ifboolKV[DS]{Calculatrice}{autorisée}{interdite} --};
    }
    ]%
    \sffamily Devoir surveillé \no\useKV[DS]{Numero} (Sujet \useKV[DS]{Sujet})\hfill\useKV[DS]{Niveau}\ieme{} \useKV[DS]{Classe}
    \par{\tiny\jobname}\hfill{\scriptsize\useKV[DS]{Date}}
  \end{tcolorbox}%
}
\renewcommand\TikzIE{%
  \begin{minipage}{0.6\linewidth}
    \begin{tcolorbox}[colback=gray!5,
      enhanced,%
      overlay unbroken and first={%
        \node[yshift=1em] at (frame.south) {\scriptsize\sffamily-- Calculatrice \ifboolKV[IE]{Calculatrice}{autorisée}{interdite} --};
      }
      ]%
      \sffamily \useKV[IE]{Nom} \no\useKV[IE]{Numero} : \useKV[IE]{Theme} \ifboolKV[IE]{Sujets}{\scriptsize(Sujet \useKV[IE]{NumSujet})}{}
      \vspace{1em}
      \par{\tiny\jobname}\hfill{\scriptsize\useKV[DS]{Date}}
    \end{tcolorbox}%
  \end{minipage}\hfill\begin{minipage}{0.35\linewidth}
    Nom : \dotfill\\[0.5em]
    Prénom : \dotfill\\[0.5em]
    Classe : \dotfill
  \end{minipage}
  \par\bigskip
}

\renewcommand\TikzFiche{%
  \begin{tcolorbox}[%spread outwards=-1cm,spread inwards=-1cm,
    colback=gray!5]%
  \Large\sffamily \useKV[Fiche]{Theme}\hfill\useKV[Fiche]{Niveau}\ieme{} \useKV[Fiche]{Classe}
  \par{\tiny\jobname}\hfill{\scriptsize\useKV[Fiche]{Date}}
  \end{tcolorbox}%
}

\renewcommand\TikzParcours{%
  \begin{tcolorbox}[%spread outwards=-1cm,spread inwards=-1cm,
    colback=gray!5]%
    \Large\sffamily \useKV[Parcours]{Theme}\hfill\useKV[Parcours]{Niveau} \useKV[Parcours]{Classe}%
    \par{\tiny\useKV[Parcours]{Code}}\hfill{\scriptsize\useKV[Parcours]{Date}}%
  \end{tcolorbox}%
}

\renewcommand\TikzPdT{%
  \begin{tcolorbox}[%spread outwards=-1cm,spread inwards=-1cm,
    colback=gray!5]%
    \Large\sffamily \useKV[PdT]{Theme}\hfill\useKV[PdT]{Niveau} \useKV[PdT]{Classe}%
    \par{\tiny\useKV[PdT]{Code}}\hfill{\scriptsize\useKV[PdT]{Date}}%
  \end{tcolorbox}%
}
\begin{document}
\maketitle
\thispagestyle{empty}
Après la création du package \lstinline!ProfCollege!, l'idée de poursuivre la factorisation des méthodes de travail a naturellement germé. Produire des fiches, des devoirs, des évaluations\dots{} avec des méthodes différentes, cela peut vite apparaître fastidieux. Créé pour améliorer cela, cet ensemble de macro-commandes est basé, quasi exclusivement, sur les environnements \lstinline!Maquette! et \lstinline!exercice! :
\begin{itemize}
\item l'environnement \lstinline!Maquette! indique le type de document souhaité et ses caractéristiques;
\item l'environnement \lstinline!exercice! adapte la présentation et les fonctions des exercices en accord avec l'environnement \lstinline!Maquette!.
\end{itemize}
\clearpage
\section*{L'environnement \lstinline!Maquette!}
\begin{Codes}[listing only]{0.4}{}
\begin{Maquette}
      
\end{Maquette}
\end{Codes}
C'est lui qui va indiquer :
\begin{itemize}
\item le type de document souhaité (Devoir Maison / Devoir Surveillé / Interrogatin écrite / Fiche d'exercices / Parcours d'exercices / Parcours personnalisé / Plan de travail);
\item et si ce travail est corrigé (à la suite d'un exercice ou à la fin du document) ou pas;
\end{itemize}
puis il indiquera les caractéristiques du document souhaité.
\clearpage
\subsection*{Les devoirs maison}
\begin{CadreMP}
  \begin{Description}
\item[La clé \Cle{DM}]\Defaut{false}
  \begin{itemize}
  \item[La clé \Cle{Numero}]\Defaut{1}
  \item[La clé \Cle{Date}]\Defaut{\lstinline!\\today!}
  \item[La clé \Cle{Classe}]\Defaut{\{\}}
  \item[La clé \Cle{Niveau}]\Defaut{3}
  \item[La clé \Cle{Code}\footnotemark]\Defaut{\{\}}
  \end{itemize}
\end{Description}
\end{CadreMP}
\footnotetext{Outil personnel de classement.}
\begin{Codes}[]{0.4}{}
\begin{Maquette}[DM]{Numero=3,Niveau=6,Classe=Zola,Date=25 décembre 2020}
  On considère un triangle $ABC$\dots
\end{Maquette}
\end{Codes}
\clearpage
\subsection*{Les devoirs surveillés}
\begin{CadreMP}
  \begin{Description}
\item[La clé \Cle{DS}]\Defaut{false}
  \begin{itemize}
  \item[La clé \Cle{Numero}]\Defaut{1}
  \item[La clé \Cle{Date}]\Defaut{\lstinline!\\today!}
  \item[La clé \Cle{Classe}]\Defaut{}
  \item[La clé \Cle{Niveau}]\Defaut{3}
  \item[La clé \Cle{Code}]\Defaut{}
  \item[La clé \Cle{Calculatrice}] autorisée ou non\Defaut{false}
  \item[La clé \Cle{Sujet}]\Defaut{A}
  \end{itemize}
\end{Description}
\end{CadreMP}
\begin{Codes}[]{0.4}{}
\begin{Maquette}[DS]{Numero=3,Niveau=6,Classe=Zola,Date=25 décembre 2020}
  On considère un triangle $ABC$\dots
\end{Maquette}
\end{Codes}
\clearpage
\subsection*{Les évaluations écrites}
\begin{CadreMP}
  \begin{Description}
\item[La clé \Cle{IE}]\Defaut{false}
  \begin{itemize}
  \item[La clé \Cle{Numero}]\Defaut{1}
  \item[La clé \Cle{Date}]\Defaut{\lstinline!\\today!}
  \item[La clé \Cle{Classe}]\Defaut{}
  \item[La clé \Cle{Niveau}]\Defaut{3}
  \item[La clé \Cle{Code}]\Defaut{}
  \item[La clé \Cle{Calculatrice}] autorisée ou nom\Defaut{false}
  \item[La clé \Cle{Sujet}] pour de multiples sujets\Defaut{\{\}}
  \item[La clé \Cle{Theme}] Thème de l'évaluation \Defaut{\{\}}
  \item[La clé \Cle{Nom}] Pour changer le nom \og \'Evaluation\fg\Defaut{\'Evaluation}
  \end{itemize}
\end{Description}
\end{CadreMP}
\begin{Codes}[]{0.4}{}
\begin{Maquette}[IE]{Numero=3,Niveau=6,Classe=Zola,Date=25 décembre 2020,Theme=La géométrie}%    
  On considère un triangle $ABC$\dots
\end{Maquette}
\end{Codes}
\clearpage
\subsection*{Les fiches d'exercices}
\begin{CadreMP}
  \begin{Description}
\item[La clé \Cle{Fiche}]\Defaut{false}
  \begin{itemize}
  \item[La clé \Cle{Date}]\Defaut{\lstinline!\\today!}
  \item[La clé \Cle{Classe}]\Defaut{Nairobi}
  \item[La clé \Cle{Niveau}]\Defaut{6}
  \item[La clé \Cle{Theme}] Thème de la fiche \Defaut{Les nombres décimaux}
  \item[La clé \Cle{Code}]\Defaut{}
  \item[La clé \Cle{NomExercice}] pour modifier le nom des exercices\Defaut{Exercice}
  \end{itemize}
\end{Description}
\end{CadreMP}
\begin{Codes}[]{0.4}{}
\begin{Maquette}[Fiche]{Niveau=6,Classe=Zola,Date=25 décembre 2020}
  On considère un triangle $ABC$\dots
\end{Maquette}
\end{Codes}
%\begin{Maquette}[Fiche]{Niveau=6,Classe=Zola,Date=25 décembre 2020}
%  On considère un triangle $ABC$\dots
%\end{Maquette}
\clearpage
\subsection*{Les parcours}
\begin{CadreMP}
  \begin{Description}
\item[La clé \Cle{Parcours}\footnotemark]\Defaut{false}
  \begin{itemize}
  \item[La clé \Cle{Date}]\Defaut{\lstinline!\\today!}
  \item[La clé \Cle{Classe}]\Defaut{Nairobi}
  \item[La clé \Cle{Niveau}]\Defaut{6}
  \item[La clé \Cle{Theme}] Thème du parcours \Defaut{Les nombres décimaux}
  \item[La clé \Cle{Code}]\Defaut{}
  \item[La clé \Cle{NomExercice}] pour modifier le nom des exercices\Defaut{Exercice}
  \end{itemize}
\end{Description}
\end{CadreMP}
  \footnotetext{Il dispose des même clés que la clé \Cle{Fiche} et de la même présentation. Ce n'est pas là son utilité\dots}
\begin{Codes}[]{0.4}{}
\begin{Maquette}[Parcours]{Niveau=6,Classe=Zola,Date=25 décembre 2020}
  On considère un triangle $ABC$\dots
\end{Maquette}
\end{Codes}
\clearpage
\subsection*{Les parcours fléchés}
\begin{CadreMP}
  \begin{Description}
\item[La clé \Cle{PdT}\footnotemark]\Defaut{false}
  \begin{itemize}
  \item[La clé \Cle{Date}]\Defaut{\lstinline!\\today!}
  \item[La clé \Cle{Classe}]\Defaut{Nairobi}
  \item[La clé \Cle{Niveau}]\Defaut{6}
  \item[La clé \Cle{Theme}] Thème du plan de travail \Defaut{Les nombres décimaux}
  \item[La clé \Cle{Code}]\Defaut{}
  \item[La clé \Cle{NomExercice}] pour modifier le nom des exercices\Defaut{Exercice}
  \end{itemize}
\end{Description}
\end{CadreMP}
  \footnotetext{Il dispose des même clés que la clé \Cle{Fiche} et de la même présentation. Ce n'est pas là son utilité\dots}
\begin{Codes}[]{0.4}{}
\begin{Maquette}[PdT]{Niveau=6,Classe=Zola,Date=25 décembre 2020}
  On considère un triangle $ABC$\dots
\end{Maquette}
\end{Codes}
%\begin{Maquette}[PdT]{Niveau=6,Classe=Zola,Date=25 décembre 2020}
%  On considère un triangle $ABC$\dots
%\end{Maquette}
\clearpage
\subsection*{Les parcours personnalisés}
\begin{CadreMP}
  \begin{Description}
\item[La clé \Cle{ParcoursPerso}]\Defaut{false}
  \begin{itemize}
  \item[La clé \Cle{Fichier}]\Defaut{}\par est le nom du fichier {\ttfamily csv} à utiliser pour créer les parcours personnalisés.
  \end{itemize}
\end{Description}
\end{CadreMP}
\begin{Codes}[listing only]{0.4}{}
\begin{Maquette}[ParcoursPerso]{Fichier=Eleves}
  On considère un triangle $ABC$\dots
\end{Maquette}
\end{Codes}
\clearpage
Comme on peut le voir, chaque type de travail est associé à une présentation. Ces présentations sont créées par les commandes\footnote{Qu'on peut donc redéfinir. Par exemple, c'est ce que j'ai fait dans ce document pour la commande \lstinline!\\TikzFiche! : elle ne commence plus une nouvelle page.} : \lstinline!\TikzDM/\TikzDMCor! pour la clé \Cle{DM}; \lstinline!\TikzDS/\TikzDSCor! pour la clé \Cle{DS}; \lstinline!\TikzIE/\TikzIECor! pour la clé \Cle{IE}; \lstinline!\TikzFiche/\TikzFicheCor! pour la clé \Cle{Fiche}, \lstinline!\TikzParcours/\TikzParcoursCor! pour la clé \Cle{Parcours} et \lstinline!\TikzPdT! pour la clé \Cle{PdT}.

On dispose également de la commande \lstinline!\Competences! permettant de construire un tableau de\dots{} compétences données par l'utilisateur; les compétences étant séparées par le symbole \lstinline!/!.
\begin{Codes}[]{0.4}{}
\Competences{Utiliser le compas/Utiliser l'équerre}
\end{Codes}
Les notations \lstinline!NA!, \lstinline!ECA! et \lstinline!A! peuvent être redéfinies par les commandes \lstinline!\PfMCompNA!, \lstinline!\PfMCompECA! et \lstinline!\PfMCompA!.
\begin{Codes}[]{0.4}{}
\renewcommand\PfMCompNA{\rule{0pt}{3ex}\RKangry}
\renewcommand\PfMCompECA{\RKsad}
\renewcommand\PfMCompA{\RKbigsmile}
\Competences{Utiliser le compas/Utiliser l'équerre}  
\end{Codes}
Si l'on souhaite davantage de niveaux d'évaluations (ou moins), on utilisera les codes suivants :
\begin{Codes}[]{0.4}{}
\Competences[4]{0§1§2§3/Utiliser le compas/Utiliser l'équerre}
\end{Codes}
\begin{Codes}[]{0.4}{}
\Competences[2]{A§NA/Utiliser le compas/Utiliser l'équerre}
\end{Codes}
\clearpage
\section*{L'environnement \lstinline!exercice!}
\begin{Codes}[listing only]{0.4}{}
\begin{exercice}
      
\end{exercice}
\end{Codes}
\subsection*{Habillage des exercices}
L'envionnement \lstinline!exercice! doit {\em impérativement} être inclus dans un environnement \lstinline!Maquette! car le choix de document influence la présentation des exercices.
\begin{Codes}[]{0.5}{}
\begin{Maquette}[IE]{Theme=Les fonctions,Niveau=3,Classe=Gide}
  \begin{exercice}
    La fonction $f:x\mapsto3x+2$ est-elle une fonction affine ? Justifier.
  \end{exercice}
\end{Maquette}
\end{Codes}
\begin{Codes}[listing side text]{0.6}{}
\begin{Maquette}[DM]{Numero=50,Niveau=3,Classe=Gide}
  \begin{exercice}
    La fonction $f:x\mapsto3x+2$ est-elle une fonction affine ? Justifier.
  \end{exercice}
\end{Maquette}
\end{Codes}
\begin{Codes}[listing side text]{0.6}{}
\begin{Maquette}[DS]{Numero=50,Niveau=3,Classe=Gide}
  \begin{exercice}
    La fonction $f:x\mapsto3x+2$ est-elle une fonction affine ? Justifier.
  \end{exercice}
\end{Maquette}
\end{Codes}
\begin{Codes}[]{0.4}{}
\begin{Maquette}[Fiche]{Theme=Les fonctions,Niveau=3,Classe=Gide}
  \begin{exercice}
    La fonction $f:x\mapsto3x+2$ est-elle une fonction affine ? Justifier.
  \end{exercice}
\end{Maquette}
\end{Codes}
%\begin{Maquette}[Fiche]{Theme=Les fonctions,Niveau=3,Classe=Gide}
%  \begin{exercice}
%    La fonction $f:x\mapsto3x+2$ est-elle une fonction affine ? Justifier.
%  \end{exercice}
%\end{Maquette}
\subsection*{Le barème des exercices}
On remarque que certains exercices sont associés à un total de point. C'est une des clés disponibles pour l'environnement \lstinline!exercice!. Voici celles pour les barèmes :
\begin{Description}
\item[La clé \Cle{BaremeTotal}] qui affichera, dans le coin supérieur droit, le total de points de l'exercice. La valeur de cette clé :
  \begin{itemize}
  \item est {\em fixée} à {\sffamily false} pour la clé \Cle{Fiche};
  \item est positionnée à {\sffamily true} mais {\em modifiable} pour les clés \Cle{DM}, \Cle{DS} et \Cle{IE}.
  \end{itemize}
\item[La clé \Cle{BaremeDetaille}]\Defaut{false}\par qui a le même fonctionnement que la clé \Cle{BaremeTotal}.
  La commande \lstinline!\brm{}! permet la construction du barème (détaillé et total).
\item[La clé \Cle{MotPoint}]\Defaut{point}\par donnée sous forme de texte puisque le pluriel est géré.
  \begin{Codes}[listing only]{0.4}{}
\begin{Maquette}[DS]{Numero=3,Classe=Euler,Niveau=4}
  \begin{exercice}  % ici le barème est total, pas de détail : comportement par défaut     
    On considère les expressions $A=2x(3x+5)$ et $B=x(7x-1)$.
    \begin{enumerate}
    \item\brm{1} Développer l'expression $A$.
    \item\brm{1.5} Développer l'expression $B$.
    \end{enumerate}
  \end{exercice}
  \begin{exercice}[BaremeDetaille]  %ici le barème est total ET détaillé
    On considère les expressions $A=2x(3x+5)$ et $B=x(7x-1)$.
    \begin{enumerate}
    \item\brm{1} Développer l'expression $A$.
    \item\brm{1.5} Développer l'expression $B$.
    \end{enumerate}
  \end{exercice}
\end{Maquette}
\end{Codes}
\begin{Maquette}[DS]{Numero=3,Classe=Euler,Niveau=4}
  \begin{exercice}  % ici le barème est total, pas de détail : comportement par défaut     
    On considère les expressions $A=2x(3x+5)$ et $B=x(7x-1)$.
    \begin{enumerate}
    \item\brm{1} Développer l'expression $A$.
    \item\brm{1.5} Développer l'expression $B$.
    \end{enumerate}
  \end{exercice}
  \begin{exercice}[BaremeDetaille]  %ici le barème est total ET détaillé
    On considère les expressions $A=2x(3x+5)$ et $B=x(7x-1)$.
    \begin{enumerate}
    \item\brm{1} Développer l'expression $A$.
    \item\brm{1.5} Développer l'expression $B$.
    \end{enumerate}
  \end{exercice}
\end{Maquette}
\end{Description}
\clearpage
\subsection*{Source, compétence et titre}
Si on souhaite citer la source d'un exercice ou lui donner un titre, on dispose des clés :
\begin{Description}
\item[La clé \Cle{Source}] pour citer la source de l'exercice.\Defaut{\{\}}
\item[La clé \Cle{Titre}] pour nommer un exercice.\Defaut{\{\}}
\item[La clé \Cle{Competence}] pour indiquer une compétence associée à l'exercice.\Defaut{\{\}}
\begin{Codes}[listing only]{0.4}{}
\begin{Maquette}[Fiche]{Theme=Algorithmique}
  \colorlet{PfMColCpt}{Crimson}
  \colorlet{PfMColSrc}{NavyBlue}
  \begin{exercice}[Source=Olympiades 2019,Titre=Modifier des mots,Competence=Raisonner]
    Dans ce problème, on appellera {\em mot} toute suite de lettres formée des lettres A, D et G. Par exemple : ADD, A, AAADG sont des {\em mots}.
    \\Astrid possède un logiciel qui fonctionne de la manière suivante : un utilisateur entre un {\em mot} et, après un clic sur EXÉCUTER, chaque lettre A du {\em mot} (s'il y en a) est remplacée par le {\em mot} AGADADAGA. Ceci donne un nouveau {\em mot}.\\Par exemple, si l'utilisateur rentre le {\em mot} AGA, on obtient le {\em mot} AGADADAGAGAGADADAGA. Un deuxième clic sur EXÉCUTER réitère la transformation décrite ci-dessus au nouveau {\em mot}, et ainsi de suite.
    \begin{enumerate}
    \item Quels sont les {\em mots} qui restent inchangés quand on clique sur EXÉCUTER ?
    \end{enumerate}
\end{exercice}
\end{Maquette}
\end{Codes}
%\clearpage
\begin{Maquette}[Fiche]{Theme=Algorithmique}
  \colorlet{PfMColCpt}{Crimson}
  \colorlet{PfMColSrc}{NavyBlue}
  \begin{exercice}[Source=Olympiades 2019,Titre=Modifier des mots,Competence=Raisonner]
    Dans ce problème, on appellera {\em mot} toute suite de lettres formée des lettres A, D et G. Par exemple : ADD, A, AAADG sont des {\em mots}.
    \\Astrid possède un logiciel qui fonctionne de la manière suivante : un utilisateur entre un {\em mot} et, après un clic sur EXÉCUTER, chaque lettre A du {\em mot} (s'il y en a) est remplacée par le {\em mot} AGADADAGA. Ceci donne un nouveau {\em mot}.\\Par exemple, si l'utilisateur rentre le {\em mot} AGA, on obtient le {\em mot} AGADADAGAGAGADADAGA. Un deuxième clic sur EXÉCUTER réitère la transformation décrite ci-dessus au nouveau {\em mot}, et ainsi de suite.
    \begin{enumerate}
    \item Quels sont les {\em mots} qui restent inchangés quand on clique sur EXÉCUTER ?
    \end{enumerate}
\end{exercice}
\end{Maquette}
\end{Description}
\clearpage
\subsection*{Des logos ?}
Les clés décrites dans cette partie ne sont disponibles que pour les maquettes
\lstinline!Fiche! / \lstinline!Parcours!.
\begin{Description}
\item[La clé \Cle{Oral}]\Defaut{false}\par
  pour indiquer que l'exercice se résoud oralement.
\item[La clé \Cle{Calculatrice}]\Defaut{true}\par
  pour afficher que la calculatrice {\em est interdite}.
\end{Description}
  \begin{Codes}[listing only,listing options={frame=,escapechar=!}]{0.4}{}
\begin{Maquette}[Fiche]{Theme=Calcul mental}
  \begin{exercice}[Oral]
    $1+1=?$
  \end{exercice}
  \begin{exercice}[!\color{DarkGreen}\ttfamily Calculatrice!=false]
    $1+1=?$
  \end{exercice}
  % On peut mélanger les deux logos.
  \begin{exercice}[!\color{DarkGreen}\ttfamily Calculatrice!=false,Oral]
    $1+1=?$
  \end{exercice}
\end{Maquette}
\end{Codes}
\begin{Maquette}[Fiche]{Theme=Calcul mental}
  \begin{exercice}[Oral]
    $1+1=?$
  \end{exercice}
  \begin{exercice}[Calculatrice=false]
    $1+1=?$
  \end{exercice}
  \begin{exercice}[Calculatrice=false,Oral]
    $1+1=?$
  \end{exercice}
\end{Maquette}
\clearpage
\subsection*{Focus sur la maquette \lstinline!Parcours!}
Dans le cas de cette maquette, l'environnement \lstinline!exercice! dispose de la clé \Cle{Trajet} permettant de construire automatiquement le ou les parcours. La commande \lstinline!\AfficheParcours{}! représente un schéma associé à un parcours. Ce schéma se construit seul, automatiquement, après que l'enseignant ait utilisé la clé \Cle{Trajet} dans la création de ses exercices. Il faudra néanmoins deux compilations. \`A noter que ce schéma, avec l'utilisation du package \lstinline!hyperref!, dispose de liens cliquables sur les exercices sélectionnés pour le parcours.
\begin{Codes}[listing only,listing options={frame=,escapechar=!}]{0.4}{\small}
\begin{Maquette}[Parcours]{Theme=Calcul mental}
  \begin{description}
  \item[!Parcours! Padawan] \AfficheParcours{Padawan}
  \item[!Parcours! Jedï] \AfficheParcours{Jedi}
  \item[!Parcours! Grand Maître] \AfficheParcours{GrandMaitre}
  \end{description}
  \begin{exercice}[Trajet={Padawan,Jedi,GrandMaitre}]
    
  \end{exercice}
  \begin{exercice}[Trajet={Padawan,Jedi}]
    
  \end{exercice}
  \begin{exercice}[Trajet={Padawan}]
    
  \end{exercice}
\end{Maquette}
\end{Codes}
\begin{Maquette}[Parcours]{Theme=Calcul mental}
  \begin{description}
  \item[Parcours Padawan] \AfficheParcours{Padawan}
  \item[Parcours Jedï] \AfficheParcours{Jedi}
  \item[Parcours Grand Maître] \AfficheParcours{GrandMaitre}
  \end{description}
  \begin{exercice}[Trajet={Padawan,Jedi,GrandMaitre}]
    
  \end{exercice}
  \begin{exercice}[Trajet={Padawan,Jedi}]
    
  \end{exercice}
  \begin{exercice}[Trajet={Padawan}]
    
  \end{exercice}
\end{Maquette}
\clearpage
\subsection*{Focus sur la maquette \lstinline!ParcoursPerso!}
Dans le cas de cette maquette, l'environnement \lstinline!exercice! n'admet aucune clé. En effet, cette maquette utilise :
\begin{itemize}
\item et un fichier {\ttfamily *.csv} de la forme :
  \begin{Codes}[listing only]{0.4}{}
Anne,Bec,{1,2}
Paul,Isse,{1,3}
Jean,Némar,{1,4}
  \end{Codes}
\item une liste d'exercices pour produire les trois pages suivantes à l'aide du code :
\begin{Codes}[listing only]{0.4}{}
\begin{Maquette}[ParcoursPerso]{Type=Fiche,Fichier=Documentation,Theme=L'alphabet,Niveau=6,Classe=Alpha}
  \begin{exercice}
    A
  \end{exercice}
  \begin{exercice}
    B
  \end{exercice}
  \begin{exercice}
    C
  \end{exercice}
  \begin{exercice}
    D
  \end{exercice}
\end{Maquette}
\end{Codes}
\end{itemize}
\clearpage
\includepdf[pages=-]{TestParcoursPerso.pdf}
\clearpage
\subsection*{Focus sur la maquette \lstinline!PdT!}
Dans le cas de cette maquette, la présentation des exercices est modifiée : il n'y a plus de titre afin de gagner de la place. De plus, chaque exercice définit, en fonction de son numéro (non écrit mais présent), huit points d'ancrages :\lstinline!N-5!, \lstinline!S-5!, \lstinline!O-5!, \lstinline!E-5!, \lstinline!NO-5!, \lstinline!NE-5!, \lstinline!SO-5! et \lstinline!SE-5! pour respectivement les points Nord, Sud, Ouest, Est, Nord-Ouest, Nord-Est, Sud-Ouest et Sud-Est du cadre de l'exercice 5.

On pourra aussi utiliser (avec les même n\oe uds précédés de \lstinline!SOS!) l'environnement \lstinline!SOS! servant d'aide aux élèves en difficulté.

De plus, deux commandes font leur apparition :
\begin{itemize}
\item \lstinline!\CheminVrai! listant les liaisons entre exercices sous la forme \lstinline!noeud départ/noeud arrivée!;
\item \lstinline!\CheminFaux! listant les liaisons entre exercices sous la forme \lstinline!noeud départ/noeud arrivée!.
\end{itemize}
L'exemple suivant est obtenu par le code présenté à la page suivante.
\begin{Maquette}[PdT]{Theme=Plan de travail : Calcul littéral,Niveau=3,Classe=Alpha,Date={},Code={}}
  \setcounter{PfMExo}{0}
    \begin{minipage}{0.35\linewidth}
      \begin{exercice}%Exercice 1
        A
      \end{exercice}
    \end{minipage}
    \hfill
    \begin{minipage}{0.3\linewidth}
      \begin{exercice}%Exercice 2
        B
      \end{exercice}
    \end{minipage}
    \hfill
    \begin{minipage}{0.15\linewidth}
      \begin{SOS}%SOS 1
        Aide A
      \end{SOS}
    \end{minipage}
    
    \vspace*{2cm}
    
    \begin{minipage}{0.35\linewidth}
      \begin{exercice}%Exercice 3
        C
      \end{exercice}
    \end{minipage}
    \hfill
    \begin{minipage}{0.3\linewidth}
      \begin{exercice}%Exercice 4
        D
      \end{exercice}
    \end{minipage}
    \hfill
    \begin{minipage}{0.15\linewidth}
      \begin{SOS}%SOS 2
        Aide B
      \end{SOS}
    \end{minipage}
    \CheminVrai{S-1/NO-3}
    \CheminFaux{E-1/O-2,E-2/SOS-O-1}
  \end{Maquette}
  \clearpage
\begin{Codes}[listing only]{}{}
  \begin{Maquette}[PdT]{Theme=Plan de travail : Calcul littéral,Niveau=3,Classe=Alpha,Date={},Code={}}
    \begin{minipage}{0.35\linewidth}
      \begin{exercice}%Exercice 1
        A
      \end{exercice}
    \end{minipage}
    \hfill
    \begin{minipage}{0.3\linewidth}
      \begin{exercice}%Exercice 2
        B
      \end{exercice}
    \end{minipage}
    \hfill
    \begin{minipage}{0.15\linewidth}
      \begin{SOS}%SOS 1
        Aide A
      \end{SOS}
    \end{minipage}
    
    \vspace*{2cm}
    
    \begin{minipage}{0.35\linewidth}
      \begin{exercice}%Exercice 3
        C
      \end{exercice}
    \end{minipage}
    \hfill
    \begin{minipage}{0.3\linewidth}
      \begin{exercice}%Exercice 4
        D
      \end{exercice}
    \end{minipage}
    \hfill
    \begin{minipage}{0.15\linewidth}
      \begin{SOS}%SOS 2
        Aide B
      \end{SOS}
    \end{minipage}
    \CheminVrai{S-1/NO-3}
    \CheminFaux{E-1/O-2,E-2/SOS-O-1}
  \end{Maquette}
\end{Codes}
\clearpage
\subsection*{\og Correction\fg{} des exercices}
  Un exercice avec correction aura la forme suivante :
  \begin{Codes}[listing only]{0.4}{}
\begin{exercice}
    
\end{exercice}
\begin{Solution}
  
\end{Solution}
\end{Codes}
On peut gérer la correction des exercices à deux niveaux :
\begin{Description}
  \item[Au niveau de l'environnement ]\lstinline!Maquette!, on dispose de deux clés :
\end{Description}
  \begin{Description}
\item[la Clé \Cle{CorrigeApres}]\Defaut{false}\par qui affiche le corrigé, s'il existe, après l'énoncé de l'exercice.
\item[la Clé \Cle{CorrigeFin}]\Defaut{false}\par qui affiche les corrigés, s'ils existent, à la fin du document.
\end{Description}

Pour la maquette \lstinline!ParcoursPerso!, seule la clé \Cle{CorrigeFin} est autorisée.

\begin{Codes}[listing only]{0.4}{\footnotesize}
\begin{Maquette}[IE,CorrigeApres]{Theme=Calcul littéral}
  \begin{exercice}% n'a pas de correction
    \begin{enumerate}
    \item $1+1=?$.
    \end{enumerate}
  \end{exercice}
  \begin{exercice}
    \begin{enumerate}
    \item Développer $A=2(x+3)$.
    \end{enumerate}
  \end{exercice}
  \begin{Solution}
    \begin{enumerate}
    \item $A=2(x+3)=2\times x+2\times3=2x+6$
    \end{enumerate}
  \end{Solution}
\end{Maquette}
\end{Codes}
\begin{Maquette}[IE,CorrigeApres]{Theme=Calcul littéral}
    \begin{exercice}% n'a pas de correction
    \begin{enumerate}
    \item $1+1=?$.
    \end{enumerate}
  \end{exercice}
  \begin{exercice}
    \begin{enumerate}
    \item Développer $A=2(x+3)$.
    \end{enumerate}
  \end{exercice}
  \begin{Solution}
    \begin{enumerate}
    \item $A=2(x+3)=2\times x+2\times3=2x+6$
    \end{enumerate}
  \end{Solution}
\end{Maquette}
\begin{Description}
\item[Au niveau de l'environnement ]\lstinline!exercice!, on peut affiner la correction proposée en écrivant le code de l'exercice sous la forme :
\end{Description}
\begin{Codes}[listing only]{0.4}{}
\begin{exercice}
    
\end{exercice}
\begin{Solution}
  
\end{Solution}
\begin{Reponse}
  
\end{Reponse}
\begin{Indice}
  
\end{Indice}
\end{Codes}
L'environnement \lstinline!exercice! dispose alors des clés suivantes :
\begin{Description}
\item[la Clé \Cle{PasCorrige}]\Defaut{false}\par supprime l'affichage de tous types de correction pour l'exercice considéré;
\item[la Clé \Cle{Pouce}]\Defaut{false}\par qui affiche uniquement le contenu de l'environnement \lstinline!Indice!;
\item[la Clé \Cle{Direct}]\Defaut{false}\par qui affiche uniquement le contenu de l'environnement \lstinline!Reponse!.
\end{Description}

\bigskip

Cela engendre, si nécessaire, la création de fichier \lstinline!*.sol! pour les solutions, \lstinline!*.rep! pour les réponses et \lstinline!*.cdp! pour les indices. Il conviendra de les effacer à chaque changement de clé (\Cle{PasCorrige}, \Cle{Pouce}, \Cle{Direct}) ou ajout d'exercice(s). En effet, la priorité d'affichage est donnée aux fichiers \lstinline!*.sol! puis \lstinline!*.rep! et enfin \lstinline!*.cdp!.

\bigskip

\begin{Codes}[listing only]{0.4}{}
\begin{Maquette}[DS,CorrigeFin]{Numero=3}
  \begin{exercice} % La clé CorrigeFin utilise le contenu de l'environnement Solution
    $1+\dfrac2{10}=$
  \end{exercice}
  \begin{Solution}
    $1+\dfrac2{10}=\num{1.2}$
  \end{Solution}
  \begin{Reponse}
    \num{1.2}
  \end{Reponse}
  \begin{Indice}
    Une unité représente dix dixièmes de l'unité.
  \end{Indice}
  \begin{exercice}[PasCorrige] % Malgré les environnements Solution, Reponse, Indice
      $7(x+5)=$ ?
  \end{exercice}
  \begin{Solution}
    $7(x+5)=7\times x+7\times 5=7x+35$
  \end{Solution}
  \begin{Indice}
    Utiliser la simple distributivité.
  \end{Indice}
  \begin{Reponse}
    $7x+35$
  \end{Reponse}
  \begin{exercice}[Direct] % On affiche uniquement l'environnement Reponse
    $\dfrac34-\dfrac25=$?
  \end{exercice}
  \begin{Solution}
    $\dfrac34-\dfrac25=\dfrac{15}{20}-\dfrac8{20}=\dfrac7{20}$
  \end{Solution}
  \begin{Reponse}
    $\dfrac7{20}$
  \end{Reponse}
  \begin{Indice}
    C'est une soustraction, il faut réduire au même dénominateur.
  \end{Indice}
  \begin{exercice}[Pouce] % On affiche uniquement l'environnement Indice
    Si $x=3$, que vaut $x^3+4x$ ?
  \end{exercice}
  \begin{Solution}
    $x^3+4x=3^3+4\times3=27+12=39$
  \end{Solution}
  \begin{Indice}
    Que veut dire $x^3$ ? Que veut dire $4x$ ?
  \end{Indice}
  \begin{Reponse}
    39
  \end{Reponse}
\end{Maquette}
\end{Codes}
\begin{Maquette}[DS,CorrigeFin]{Numero=3}
  \begin{exercice} % La clé CorrigeFin utilise le contenu de l'environnement Solution
    $1+\dfrac2{10}=$
  \end{exercice}
  \begin{Solution}
    $1+\dfrac2{10}=\num{1.2}$
  \end{Solution}
  \begin{Reponse}
    \num{1.2}
  \end{Reponse}
  \begin{Indice}
    Une unité représente dix dixièmes de l'unité.
  \end{Indice}
  \begin{exercice}[PasCorrige]
    $7(x+5)=$ ?
  \end{exercice}
  \begin{Solution}
    $7(x+5)=7\times x+7\times 5=7x+35$
  \end{Solution}
  \begin{Indice}
    Utiliser la simple distributivité.
  \end{Indice}
  \begin{Reponse}
    $7x+35$
  \end{Reponse}
  \begin{exercice}[Direct]
    
    $\dfrac34-\dfrac25=$?
  \end{exercice}
  \begin{Solution}
    $\dfrac34-\dfrac25=\dfrac{15}{20}-\dfrac8{20}=\dfrac7{20}$
  \end{Solution}
  \begin{Reponse}
    $\dfrac7{20}$
  \end{Reponse}
  \begin{Indice}
    C'est une soustraction, il faut réduire au même dénominateur.
  \end{Indice}
  \begin{exercice}[Pouce]
    Si $x=3$, que vaut $x^3+4x$ ?
  \end{exercice}
  \begin{Solution}
    $x^3+4x=3^3+4\times3=27+12=39$
  \end{Solution}
  \begin{Indice}
    Que veut dire $x^3$ ? Que veut dire $4x$ ?
  \end{Indice}
  \begin{Reponse}
    39
  \end{Reponse}
\end{Maquette}
\end{document}
%%% Local Variables: 
%%% TeX-engine: luatex
%%% End: